\section{The Birth of Human Language}
\label{sec:birth-of-language}
\slidesrc{Intro: 4--21}

Before teaching computers to process language, it is worth asking what language
\emph{is} and where it comes from. This section is not technical: it frames the object
of study of the whole course. The questions it raises (what makes human language
special? why did it appear?) come back later, for example when we see that language is
used not only to describe the world but also to \emph{act} on other people
(pragmatics, Section~\ref{sec:linguistic-analysis}).

This section covers:
\begin{enumerate}[nosep]
  \item what distinguishes human language from animal communication
        (Section~\ref{sec:human-vs-animal});
  \item when language appeared, on the timeline of human evolution
        (Section~\ref{sec:timeline});
  \item the evolutionary pressures that may have pushed our species towards language,
        and the main theories about how it emerged (Section~\ref{sec:drivers});
  \item the ``gossip'' hypothesis: language as a tool for managing social relationships
        (Section~\ref{sec:gossip}).
\end{enumerate}

\subsection{Human Language vs.\ Animal Communication}
\label{sec:human-vs-animal}
\slidesrc{Intro: 5--8}

Many animals communicate: birds sing to mark territory, vervet monkeys have different
alarm calls for eagles, leopards and snakes, dogs growl and wag their tails. So what is
special about \emph{human} language? The course poses this as an open question and
gives two hints.

\subsubsection{Hint 1: talking about what is not here and now}
\slidesrc{Intro: 5--7}

The first hint is given only through two pictures: a portrait of the Latin poet
\textbf{Horace} and the poster of the film \textbf{\emph{Dead Poets Society}}. What links
them is the phrase \emph{carpe diem} (``seize the day''): it comes from one of Horace's
odes, and it is the motto that the teacher in the film gives his students.

The most likely reading of the hint is this: \emph{carpe diem} is a sentence about
\emph{time}. It urges you to act now because the future is uncertain and the past is
gone. To say it at all, a language must be able to talk about things that are not in front
of the speaker: the future, the past, possibilities, abstract ideas. Animal signals almost
never do this: an alarm call means ``danger \emph{here}, \emph{now}''.

\begin{nbdef}{Displacement}
A communication system has \textbf{displacement} if its messages can refer to things
that are \emph{removed in time or space} from the moment and place of communication:
past and future events, distant places, hypothetical or imaginary entities, abstract
concepts.
\end{nbdef}

Displacement is one of the ``design features'' of language listed by the linguist
Charles Hockett (1960). Others worth knowing, because they will come back in this course,
are:
\begin{itemize}[nosep]
  \item \textbf{productivity}: a finite set of words and rules generates an unbounded
        number of new sentences, never heard before, that are still understood;
  \item \textbf{duality of patterning}: meaningless units (sounds, letters) combine into
        meaningful units (morphemes, words; see Section~\ref{sec:morphology}), which
        combine into sentences;
  \item \textbf{arbitrariness}: there is no natural link between a word and what it
        means (\emph{dog}, \emph{cane}, \emph{Hund} all name the same animal);
  \item \textbf{cultural transmission}: a language is learnt from the community, not
        inherited genetically.
\end{itemize}

\subsubsection{Hint 2: the bees, an exception}
\slidesrc{Intro: 8}

The second hint says that the language of bees is an \emph{exception} to the most
evident difference between human language and animal communication: bees convey a
specific piece of information that no other animal conveys in the wild.

That piece of information is \emph{the location of a distant food source}. When a
forager bee comes back to the hive, it performs the \textbf{waggle dance} (decoded by Karl
von Frisch, Nobel Prize 1973):
\begin{itemize}[nosep]
  \item the \textbf{angle} $\alpha$ of the straight ``waggle run'' with respect to the
        vertical of the honeycomb equals the angle between the direction of the food and
        the direction of the sun;
  \item the \textbf{duration} of the waggle run encodes the \textbf{distance} to the food
        (longer run, farther food).
\end{itemize}

\begin{center}
\begin{tikzpicture}[>=Stealth, scale=1.05]
  % left panel: outside the hive
  \begin{scope}
    \node[font=\small\bfseries] at (1.5,4.3) {In the field};
    \fill[nborange!80] (0,0) circle (0.12) node[below=3pt, black, font=\scriptsize] {hive};
    \draw[->, thick, nbgray] (0,0) -- (0,3.4) node[pos=0.85, left, font=\scriptsize, black] {towards the sun};
    \draw[nborange, fill=nborange!30] (0,3.75) circle (0.18);
    \draw[->, very thick, nbred] (0,0) -- (2.2,2.6);
    \node[nbred, font=\scriptsize, right] at (2.2,2.6) {food};
    \draw (0,1.2) arc[start angle=90, end angle=49.8, radius=1.2];
    \node at (0.42,1.45) {$\alpha$};
  \end{scope}
  % right panel: on the comb
  \begin{scope}[xshift=7cm]
    \node[font=\small\bfseries] at (0,4.3) {On the vertical comb};
    \draw[->, thick, nbgray] (0,0.4) -- (0,3.5) node[above, font=\scriptsize, black] {``up'' = sun};
    \draw[->, very thick, nbred] (-0.5,0.6) -- (1.3,2.7) node[above right, font=\scriptsize, align=left, xshift=2mm] {waggle run\\(duration $\propto$ distance)};
    \draw[dashed, nbgray] (1.3,2.7) .. controls (2.6,2.2) and (1.5,0.2) .. (-0.5,0.6);
    \draw[dashed, nbgray] (1.3,2.7) .. controls (-0.6,3.6) and (-1.8,1.4) .. (-0.5,0.6);
    \draw (0,1.6) arc[start angle=90, end angle=49.4, radius=1.0];
    \node at (0.35,1.95) {$\alpha$};
  \end{scope}
\end{tikzpicture}
\captionof{figure}{The waggle dance. The same angle $\alpha$ appears twice: in the field
it separates the sun from the food, on the comb it separates ``up'' from the waggle run.
The dashed loops are the return paths the bee walks between two waggle runs.}
\end{center}

So bees can talk about something \emph{far away}: their communication has a (limited)
form of spatial displacement. It is still very far from human language: it covers only
one topic (food and nest sites), it is not productive (bees cannot combine dances into
new messages about new topics) and it is innate, not learnt.

\begin{nbintuition}[the answer in one line]
Human language is (among other things) \emph{displaced}: it lets us talk about the
past, the future, the absent and the imaginary. Animal communication is almost always
about the here and now; the bees' dance is the classic exception, because it describes
a place the bees are not currently in.
\end{nbintuition}

\subsection{When Did Language Appear?}
\label{sec:timeline}
\slidesrc{Intro: 9--12}

\subsubsection{The timeline of humankind}
\slidesrc{Intro: 9--10}

Table~\ref{tab:harari} places our species on the timeline of human evolution. It is
adapted from Yuval Noah Harari's \emph{Sapiens: A Brief History of Humankind} (2014).
The abbreviations are MYA = million years ago, YA = years ago.

\begin{center}
\small
\begin{tabularx}{\linewidth}{@{}l X@{}}
\toprule
\textbf{When} & \textbf{What happened} \\
\midrule
6 MYA            & Last common ``grandmother'' of humans and chimpanzees \\
2.5 MYA          & Evolution of the genus \emph{Homo} in Africa; first stone tools \\
2 MYA            & \textbf{Humans spread from Africa to Eurasia}; evolution of different human
                   species (\emph{Homo rudolfensis} and \emph{Homo ergaster} in Africa,
                   \emph{Homo erectus} in East Asia) \\
500{,}000 YA     & Neanderthals evolve in Europe and the Middle East \\
200{,}000 YA     & \emph{Homo sapiens} evolves in East Africa \\
70{,}000 YA      & \textbf{Sapiens spreads out of Africa} \\
70{,}000--30{,}000 YA & \textbf{Cognitive revolution}: invention of boats, arrows, needles;
                   beginnings of art, religion, commerce, social stratification,
                   \textbf{and language} \\
45{,}000 YA      & Sapiens settles Australia \\
30{,}000 YA      & Extinction of Neanderthals \\
16{,}000 YA      & Sapiens settles America \\
13{,}000 YA      & Extinction of \emph{Homo floresiensis}: \textbf{Homo sapiens is the only
                   surviving human species} \\
12{,}000 YA      & Agricultural revolution, permanent settlements \\
\bottomrule
\end{tabularx}
\captionof{table}{Milestones of human evolution (adapted from Harari). Note how recent
everything after the cognitive revolution is compared with the 6 million years of the
whole story.}
\label{tab:harari}
\end{center}

The course also shows a map of the movement of \emph{Homo sapiens} out of Africa: from
East Africa into the Middle East, then along the coasts of southern Asia to Australia
(about 50{,}000 years ago), into Europe and northern Asia, and finally across the Bering
region into America (about 15{,}000 years ago), while Neanderthals already lived in
Europe and \emph{Homo erectus} in Asia.

\subsubsection{What recent research says}
\slidesrc{Intro: 11}

Harari places the birth of language in the cognitive revolution, about
70{,}000--30{,}000 years ago. A 2025 study reported by MIT News (``When did human language
emerge?'', 14 March 2025) argues for an earlier date:

\begin{itemize}[nosep]
  \item \textbf{135{,}000 years ago}: genomic and geographic data suggest that the
        biological and cognitive \emph{capacity} for language was already fundamentally
        present. The argument is that the earliest split between human populations
        happened about 135{,}000 years ago, and since \emph{all} human populations have
        language, the capacity must predate the split.
  \item \textbf{100{,}000 years ago}: archaeological records show an explosion of
        symbolic behaviours, like art and complex tools, indicating active social
        communication. Language was being \emph{used} socially.
  \item \textbf{Gradual origins}: scientists widely agree that language developed from a
        complex mix of hand gestures, facial expressions and vocalisations over time,
        not in a single step.
\end{itemize}

\begin{nbwarning}[capacity vs.\ use]
The two dates answer different questions. ``The capacity for language existed
135{,}000 years ago'' is about biology (the brain and vocal tract could support
language). ``Symbolic behaviour exploded 100{,}000 years ago'' is about culture (people
were actually communicating with symbols). Do not mix them up when asked ``when did
language appear?''.
\end{nbwarning}

\subsubsection{Did the Neanderthals speak?}
\slidesrc{Intro: 12}

The course leaves this as a question for discussion. There is no definitive answer,
but these are the main pieces of evidence usually brought up:
\begin{itemize}[nosep]
  \item \textbf{for}: Neanderthals had the same variant as ours of the \emph{FOXP2} gene,
        involved in speech and language; a Neanderthal hyoid bone (the bone that
        supports the tongue) found at Kebara, Israel, is almost identical to ours; they
        made complex tools, used fire, buried their dead and probably made simple
        symbolic objects;
  \item \textbf{against / uncertain}: their symbolic record is much poorer than that of
        \emph{Sapiens} of the same period, and soft tissue (the vocal tract, the brain
        itself) does not fossilise, so the anatomy of speech can only be inferred.
\end{itemize}
Many researchers today think Neanderthals had \emph{some} form of language, though
possibly less complex than ours. Note that if language capacity predates 135{,}000
years ago in \emph{Sapiens}, the question of whether it also predates the split between
\emph{Sapiens} and Neanderthals (several hundred thousand years ago) remains open.

\subsection{Why Did Language Evolve?}
\label{sec:drivers}
\slidesrc{Intro: 13--18}

Knowing \emph{when} language appeared does not tell us \emph{why}. This part of the
course follows the linguist Sverker Johansson (\emph{The Dawn of Language: How We Came
to Talk}), for whom one major driver was the \textbf{obstetrical dilemma}.

\subsubsection{The obstetrical dilemma}
\slidesrc{Intro: 13--16}

Two evolutionary trends pulled the human body in opposite directions:
\begin{itemize}[nosep]
  \item \textbf{Upright walking (bipedalism)} restructured the human pelvis into a
        basket-like shape that supports the internal organs and makes walking efficient,
        but \emph{narrows the birth canal}.
  \item \textbf{Large brains}: hominid brain size increased dramatically over millions
        of years, so the baby's head became large compared with the mother's pelvis.
\end{itemize}
\textbf{Consequence}: labour and delivery are difficult and dangerous in the human
species (not only in humans, but especially in humans).

Evolution's ``compromise'' was to have babies born \textbf{immature}: their brain keeps
growing a lot after birth. Human newborns cannot even cling to their mother, as the
babies of all other higher primates do, and they need the protection of their family
or social group for a very long time.

\begin{nbintuition}[how this leads to language]
The course asks: how might these facts relate to the birth of language? Follow the
chain:
\begin{enumerate}[nosep]
  \item A difficult birth means a mother needs help: \textbf{midwives} and social
        support. Giving birth becomes a social, cooperative event.
  \item A baby that cannot cling must be carried or \textbf{put down}. A mother who puts
        the baby down to gather food can still soothe and monitor it from a distance with
        her \textbf{voice} (``motherese'').
  \item A long childhood is a long \textbf{window for learning}: there is time to
        absorb something as complex as a language from the community.
\end{enumerate}
\end{nbintuition}

Table~\ref{tab:pressures} summarises the argument.

\begin{center}
\small
\begin{tabularx}{\linewidth}{@{}>{\raggedright\arraybackslash}p{2.6cm}
                               >{\raggedright\arraybackslash}p{4.1cm} X@{}}
\toprule
\textbf{Evolutionary pressure} & \textbf{Physical consequence} & \textbf{Behavioural and language outcome} \\
\midrule
Bipedalism and big brains & Narrow birth canal; difficult, dangerous labour &
  Need for \textbf{midwives and social support}, forcing cooperative social structures \\
Neurological immaturity & Babies are born completely helpless and cannot cling &
  Mothers must \textbf{``put the baby down''}, using vocalisations to soothe and monitor
  them from afar \\
Extended childhood & Years of dependency requiring group protection &
  Long childhoods provide a \textbf{massive window for learning} complex language from
  the community \\
\bottomrule
\end{tabularx}
\captionof{table}{From the obstetrical dilemma to language.}
\label{tab:pressures}
\end{center}

\subsubsection{Other key drivers}
\slidesrc{Intro: 17}

Besides the obstetrical dilemma, three more drivers are usually cited:
\begin{itemize}[nosep]
  \item \textbf{Survival and cooperation}: early humans needed to coordinate group hunts,
        plan defences against predators and share crucial resources.
  \item \textbf{Social bonding}: communication grew from basic physical grooming into
        storytelling around fires, building community trust and deep empathy (this is
        developed in Section~\ref{sec:gossip}).
  \item \textbf{Abstract thought}: language enabled humans to refer to things not
        immediately present, like the past, the future or imaginary cultural concepts.
        This is displacement again (Section~\ref{sec:human-vs-animal}).
\end{itemize}

\subsubsection{Three theories of how language emerged}
\slidesrc{Intro: 18}

There are three main families of theories. They differ on \emph{speed} (gradual or
sudden) and on the \emph{starting medium} (voice or gestures).

\begin{center}
\small
\begin{tabularx}{\linewidth}{@{}>{\raggedright\arraybackslash}p{1.9cm} *{3}{>{\raggedright\arraybackslash}X}@{}}
\toprule
 & \textbf{Continuity theory} & \textbf{Discontinuity theory} & \textbf{Gesture-first theory} \\
\midrule
Core idea & Language evolved \textbf{gradually} over millions of years out of animal
  communication systems &
  Language appeared \textbf{suddenly} as a unique, exclusively human genetic mutation &
  Language originated from \textbf{physical gestures} and manual signs before shifting
  to speech \\[2pt]
Key mechanism & Natural selection slowly shaped the vocal tract and brain tissue to
  handle complex communication &
  A sudden cognitive ``leap'' or single genetic mutation re-wired the brain for
  universal grammar &
  Bipedalism freed the hands, allowing early humans to use detailed hand and facial
  signs to communicate \\[2pt]
Primary driver & The gradual need for social bonding and survival coordination in
  expanding groups &
  An internal computational upgrade in the brain, prioritising thought organisation over
  communication &
  The mechanical advantage of visual ``miming'' before vocalisation evolved to take
  over \\[2pt]
Famous proponents & Charles Darwin & Noam Chomsky & Michael Corballis \\
\bottomrule
\end{tabularx}
\captionof{table}{The three main theories on the emergence of language.}
\label{tab:theories}
\end{center}

\begin{nbintuition}[telling them apart]
\begin{itemize}[nosep]
  \item \textbf{Darwin}: language is like any other trait, shaped slowly by natural
        selection; ``more of'' what animals already have.
  \item \textbf{Chomsky}: language is primarily a tool for \emph{thinking}, and its core
        (the ability to combine units recursively, ``universal grammar'') appeared in one
        jump; communication came second.
  \item \textbf{Corballis}: language started in the \emph{hands}; the voice took over
        later because it works in the dark, at a distance and with the hands busy.
\end{itemize}
Chomsky's name will come back in this course for a very different reason: his work on
formal grammars underlies the syntactic parsing we will study later.
\end{nbintuition}

\subsection{The Role of Gossip}
\label{sec:gossip}
\slidesrc{Intro: 19--21}

One of the most popular explanations of \emph{why} language evolved is surprisingly
down-to-earth: \textbf{gossip}. The idea is due to the anthropologist Robin Dunbar
(\emph{Grooming, Gossip and the Evolution of Language}, 1996), and Harari's
\emph{Sapiens} makes it widely known.

\begin{nbfact}{Gossip as vocal grooming}
Gossip acted as a \emph{vocal substitute for physical grooming} in primates, allowing
early humans to bond and manage social relationships in groups too large for direct
touch alone.
\end{nbfact}

The argument has three steps:
\begin{enumerate}[nosep]
  \item \textbf{Primate roots.} Primates maintain social bonds and group cohesion through
        physical grooming (picking ticks and fleas from each other).
  \item \textbf{The group size limit.} Grooming is one-to-one and takes time. As early
        human groups grew larger, physical grooming became too time-consuming to bond
        with every member.
  \item \textbf{Vocal grooming.} Language evolved as a form of ``vocal grooming'' that
        lets one person bond with \emph{several} others at the same time, instead of one
        at a time.
\end{enumerate}

\begin{nbintuition}[why the numbers matter]
If grooming one partner takes a fixed amount of time per day, the time a primate must
spend grooming grows with the size of its group, and there are only so many hours in a
day. Dunbar estimated that the natural human group size (``Dunbar's number'', about 150
people) would require spending around 40\% of the day grooming. Talking to three or four
listeners at once removes that bottleneck.
\end{nbintuition}

\subsubsection{Functions of early gossip}
\slidesrc{Intro: 21}

Gossip was not idle chatter. It did three useful jobs for the group:
\begin{itemize}[nosep]
  \item \textbf{Information sharing}: it transmitted crucial \emph{reputational data}:
        who could be trusted, and who was a free rider (someone who takes advantage of
        the group without contributing).
  \item \textbf{Social cohesion}: it strengthened group identity, enforced social norms
        and managed complex social networks without written records.
  \item \textbf{Conflict reduction}: it reduced friction inside the group by aligning
        moral judgements and expectations.
\end{itemize}

\begin{nbwhy}
Two ideas from this section shape how we look at language in NLP:
\begin{itemize}[nosep]
  \item language is mostly about \emph{people} and their relationships, not just a code
        for describing facts. A lot of real text (social media, chats, reviews) is
        exactly ``gossip'' in this sense, and analysing it means dealing with opinions,
        emotions and social intentions;
  \item language is used to \emph{do} things to others (to bond, to warn, to persuade),
        which is the subject of \textbf{pragmatics}
        (Section~\ref{sec:linguistic-analysis}).
\end{itemize}
\end{nbwhy}

\begin{nbsummary}[the birth of human language]
\begin{itemize}[nosep]
  \item Key difference from animal communication: \textbf{displacement} (talking about
        the past, the future, the absent, the imaginary). The bees' waggle dance is the
        classic exception: it describes the direction and distance of distant food.
  \item Harari: language arises with the \textbf{cognitive revolution},
        70{,}000--30{,}000 years ago. Recent studies: the \emph{capacity} was present by
        \textbf{135{,}000} years ago, symbolic \emph{use} exploded around
        \textbf{100{,}000} years ago; origins were gradual (gestures, faces, voice).
  \item \textbf{Obstetrical dilemma} (Johansson): bipedalism + big brains
        $\Rightarrow$ difficult births and immature babies $\Rightarrow$ need for social
        support, vocal soothing at a distance and a long learning window.
  \item Other drivers: survival and cooperation, social bonding, abstract thought.
  \item Theories: \textbf{continuity} (gradual; Darwin), \textbf{discontinuity}
        (sudden mutation; Chomsky), \textbf{gesture-first} (hands before voice;
        Corballis).
  \item \textbf{Gossip} (Dunbar): language as vocal grooming for groups too large for
        physical grooming; it shares reputational data, builds cohesion and reduces
        conflict.
\end{itemize}
\end{nbsummary}

